\documentclass[aspectratio=169,10pt]{beamer}
\usepackage[T1]{fontenc}
\usepackage[utf8]{inputenc}
\usepackage[spanish,es-nodecimaldot]{babel}
\usepackage{lmodern}
\usepackage{amsmath}
\usepackage{booktabs}
\usepackage{graphicx}

% ============================================================
% TEMA (mismo que la presentación)
% ============================================================
\usetheme{Madrid}
\definecolor{azulfi}{RGB}{11,110,153}
\usecolortheme[named=azulfi]{structure}
\setbeamercolor{block title}{fg=black,bg=azulfi!20}
\setbeamercolor{block body}{bg=azulfi!2}
\setbeamercolor{block title alerted}{fg=white,bg=red!75!black}
\setbeamercolor{block body alerted}{bg=red!75!black!10}
\setbeamercolor{block title example}{fg=white,bg=azulfi!80!black}
\setbeamercolor{block body example}{bg=azulfi!10}
\setbeamerfont{block title}{size=\small}
\setbeamertemplate{navigation symbols}{}
\setbeamersize{text margin left=3.5mm,text margin right=3.5mm}
\graphicspath{{figuras/}}

% Pregunta (título del bloque) y respuesta (cuerpo)
\newcommand{\qa}[2]{\begin{block}{#1}\footnotesize #2\end{block}}

\title[Banco de preguntas]{Banco de preguntas para la defensa}
\subtitle{Detección asistida de melanoma en imágenes dermatoscópicas}
\institute[UACH-FING]{Universidad Autónoma de Chihuahua\\Facultad de Ingeniería}
\date{27 de noviembre de 2026}
\author[E. E. Ontiveros Lara]{Edgar Eduardo Ontiveros Lara}
\titlegraphic{\includegraphics[height=14mm]{logo_uach}\hspace{6.8cm}\includegraphics[height=14mm]{logo_fi}}

\begin{document}

\begin{frame}[plain]
  \titlepage
\end{frame}

\begin{frame}{Cómo usar este material}
  \begin{itemize}
    \item Son \textbf{47 preguntas} probables, en el orden de la presentación. Cada sección indica las láminas a las que corresponde.
    \item Las respuestas son una \textbf{guía}, no un guion: hay que poder decirlas con palabras propias.
    \item Método sugerido: tapar la respuesta, contestar en voz alta y después comparar.
    \item Todas las cifras salen de los reportes del repositorio (\texttt{reports/} y \texttt{docs/}). Si una cifra no se recuerda, es mejor decir el orden de magnitud que inventarla.
    \item Si no se sabe una respuesta: \textit{«eso no lo medí; lo que sí puedo decir es...»}.
  \end{itemize}
  \vspace{0.6em}
  \begin{exampleblock}{Las cinco cifras que hay que saber de memoria}
    Prevalencia \textbf{1.76\,\%} \quad$\cdot$\quad AUC-ROC \textbf{0.843} \quad$\cdot$\quad AUPRC \textbf{0.113} \quad$\cdot$\quad \textbf{84 de 86} melanomas detectados \quad$\cdot$\quad especificidad \textbf{39\,\%}
  \end{exampleblock}
\end{frame}

% ============================================================
\section{Problema y alcance}
\begin{frame}{Problema y alcance (láminas 3 a 6)}
  \qa{1. ¿Por qué «detección asistida» y no «diagnóstico»?}{
    El sistema entrega una probabilidad estadística sobre una imagen. No está validado clínicamente ni aprobado por una autoridad sanitaria. Su salida es «referir» o «no referir» a dermatología; la decisión sigue siendo del médico.}
  \qa{2. ¿Qué estaba mal con el 90.39\,\% de exactitud de la primera versión?}{
    Cuatro cosas: se midió sobre una validación contaminada (sin verificar fugas), con 48\,\% de melanomas cuando la prevalencia real es 1.76\,\%, con el umbral 0.5 por omisión y sin calibración. Además, la exactitud engaña cuando las clases están desbalanceadas.}
  \qa{3. ¿Por qué ISIC 2020 y no otro conjunto de datos?}{
    Porque incluye el \textbf{identificador de paciente}, que permite separar las particiones por paciente y evitar fugas. Alternativas como ISIC 2019 no lo tienen. El costo es que hay pocos melanomas: 584.}
\end{frame}

\begin{frame}{Problema y alcance (láminas 3 a 6)}
  \qa{4. ¿Por qué solo melanoma y solo dermatoscopía?}{
    Es lo que ISIC 2020 permite: sus etiquetas son melanoma contra el resto y sus imágenes son dermatoscópicas. Fuera de ese dominio el modelo no sirve, y eso se midió (DDI). Una probabilidad baja no descarta otros cánceres de piel, como el carcinoma basocelular.}
  \qa{5. ¿Se podría usar hoy en una clínica?}{
    No. Faltan una evaluación prospectiva, datos de población mexicana, validación con los equipos de la clínica y aprobación regulatoria. Es una herramienta educativa y de investigación, y además los datos son de licencia no comercial.}
  \qa{6. ¿Cuál es la aportación del trabajo, si el resultado no supera al estado del arte?}{
    La aportación es metodológica: una evaluación honesta y reproducible (particiones por paciente, conjunto de prueba bloqueado, intervalos de confianza, calibración, evaluación externa y limitaciones medidas), más un sistema completo que funciona de principio a fin.}
\end{frame}

% ============================================================
\section{Datos}
\begin{frame}{Datos (láminas 7 y 8)}
  \qa{7. ¿Qué es una fuga de datos y cómo se evitó?}{
    Es cuando información de la evaluación se cuela al entrenamiento. Aquí el riesgo era que fotos del mismo paciente quedaran en entrenamiento y en prueba. Se evitó particionando por paciente y fusionando a los pacientes que comparten imágenes duplicadas. Resultado: 0 pacientes en más de una partición, vigilado por una prueba automática.}
  \qa{8. ¿Cómo se detectaron los duplicados?}{
    En dos pasadas. Con SHA256 se encuentran archivos idénticos byte a byte: 433 grupos, todos dentro de un mismo paciente. Con un \textit{hash} perceptual (pHash de 64 bits) se encuentran imágenes visualmente iguales: 18 grupos cruzan pacientes distintos.}
  \qa{9. ¿Por qué no se eliminaron los duplicados dentro de un paciente?}{
    Porque no generan fuga: todas las imágenes de un paciente caen en la misma partición. Eliminarlos habría sido una decisión adicional sin beneficio para la validez de la evaluación.}
\end{frame}

\begin{frame}{Datos (láminas 7 y 8)}
  \qa{10. El 82\,\% de las imágenes tiene diagnóstico «unknown». ¿No es un problema?}{
    Son lesiones benignas sin subtipo; la etiqueta binaria (melanoma o no) sí está completa. La limitación real es otra: los melanomas tienen confirmación histopatológica, pero algunas benignas se etiquetaron por consenso clínico o seguimiento. Ese posible error de etiqueta tiende a \textbf{subestimar} la especificidad medida.}
  \qa{11. ¿De dónde vienen los datos y a quién representan?}{
    De seis instituciones de Europa, Estados Unidos y Australia, recopilados hasta 2020. No incluyen pacientes mexicanos, así que el desempeño en nuestra población se desconoce.}
  \qa{12. ¿Cómo se sabe que los datos descargados son los correctos?}{
    Los conteos (33,126 imágenes, 584 melanomas, 2,056 pacientes) se verificaron contra la publicación oficial, y cada imagen está registrada con su SHA256 en un manifiesto versionado.}
\end{frame}

% ============================================================
\section{Métricas}
\begin{frame}{Desbalance y métricas (lámina 9)}
  \qa{13. ¿Por qué no se reporta la exactitud?}{
    Con 1.76\,\% de melanomas, responder siempre «benigno» da 98.24\,\% de exactitud sin detectar un solo melanoma. La exactitud premia ignorar la clase que importa.}
  \qa{14. ¿Qué es el AUC-ROC?}{
    La probabilidad de que un melanoma elegido al azar reciba una puntuación mayor que una lesión benigna elegida al azar. Vale 0.5 al azar y 1.0 si el orden es perfecto. No depende del umbral.}
  \qa{15. ¿Qué es el AUPRC y por qué es la métrica primaria?}{
    El área bajo la curva de precisión contra sensibilidad. Se concentra en la clase rara: al azar vale lo mismo que la prevalencia ($\approx 0.016$). Con tantos negativos el AUC-ROC puede verse optimista; el AUPRC muestra cuántas de las alertas son melanomas de verdad.}
\end{frame}

\begin{frame}{Desbalance y métricas (láminas 9 y 16)}
  \qa{16. Defina sensibilidad, especificidad, VPP y VPN con los números del proyecto.}{
    \begin{tabular}{@{}lll@{}}
      \textbf{Sensibilidad}: de los melanomas, cuántos se refieren & 84 de 86 & 97.7\,\% \\
      \textbf{Especificidad}: de las benignas, cuántas no se refieren & 2,028 de 5,166 & 39.3\,\% \\
      \textbf{VPP}: de las referidas, cuántas son melanoma & 84 de 3,222 & 2.6\,\% \\
      \textbf{VPN}: de las no referidas, cuántas son benignas & 2,028 de 2,030 & 99.9\,\% \\
    \end{tabular}}
  \qa{17. ¿Por qué el VPP y el VPN dependen de la prevalencia y la sensibilidad no?}{
    La sensibilidad y la especificidad se calculan dentro de cada clase. El VPP y el VPN mezclan las dos clases, así que cambian con la proporción de melanomas. Proyección del modelo único: con 5\,\% de prevalencia el VPP subiría a cerca de 8\,\%.}
\end{frame}

% ============================================================
\section{Protocolo}
\begin{frame}{Protocolo de evaluación (lámina 10)}
  \qa{18. ¿Por qué el conjunto de prueba solo se abre dos veces?}{
    Cada vez que se mira la prueba y se decide algo con ella, deja de ser una estimación limpia y se convierte en otra validación. Limitar los accesos garantiza que la cifra reportada no está ajustada a esos datos.}
  \qa{19. ¿Qué es un «protocolo congelado»?}{
    Un archivo que fija el modelo, la calibración, los umbrales y las métricas, registrado en git \textbf{antes} de que exista el código que lee la prueba. El script se niega a correr si el archivo cambió, y su \textit{hash} queda anotado en la bitácora junto con la fecha del acceso.}
  \qa{20. ¿Por qué los intervalos de confianza se calculan por paciente?}{
    Las imágenes de un mismo paciente no son independientes. Remuestrear imágenes daría intervalos demasiado estrechos. Se remuestrean pacientes completos, 2,000 veces.}
\end{frame}

\begin{frame}{Protocolo de evaluación (lámina 10)}
  \qa{21. ¿Por qué estratificar por «paciente con al menos un melanoma»?}{
    Con tan pocos melanomas, una partición al azar podría dejar la prueba casi sin positivos. La estratificación reparte los pacientes con melanoma de forma proporcional; el requisito era al menos 80 melanomas en la prueba y se obtuvieron 86.}
  \qa{22. ¿Para qué sirve entonces la validación?}{
    Para todas las decisiones: línea base, tratamiento del desbalance, arquitectura, resolución, época de paro, ensamble, calibración y umbral. Por eso sus métricas están sesgadas al alza y no se reportan como resultado final.}
  \qa{23. ¿Qué pasa si ahora se quiere mejorar el modelo?}{
    El presupuesto de accesos está agotado. Cualquier mejora solo puede medirse sobre validación, o con un conjunto de prueba nuevo. Es una limitación asumida a cambio de que la cifra actual sea creíble.}
\end{frame}

% ============================================================
\section{Entrenamiento}
\begin{frame}{Entrenamiento y selección del modelo (láminas 11 y 12)}
  \qa{24. ¿Qué es el aprendizaje por transferencia y por qué se usó?}{
    Se parte de una red ya entrenada con millones de imágenes generales (ImageNet) y se ajusta a la tarea. Con solo 410 melanomas en entrenamiento, entrenar desde cero no sería viable: la red ya trae detectores de bordes, texturas y colores.}
  \qa{25. ¿Qué hace el peso a la clase positiva?}{
    Multiplica por $\approx 55$ (negativos entre positivos) la pérdida de cada melanoma: equivocarse en un melanoma cuesta lo mismo que equivocarse en 55 benignas. Efecto secundario: las probabilidades salen infladas, y por eso después se calibra.}
  \qa{26. El muestreo ponderado tuvo mejor AUPRC (0.115 contra 0.077). ¿Por qué no se eligió?}{
    Porque el sistema opera en sensibilidad alta, y ahí el muestreo deja una especificidad de 0.19 contra 0.42 del peso. Además esa diferencia de AUPRC es de una sola semilla y con intervalos traslapados.}
\end{frame}

\begin{frame}{Entrenamiento y selección del modelo (láminas 11 y 12)}
  \qa{27. ConvNeXt-T tuvo los mejores números. ¿Por qué se eligió EfficientNetV2-S a 224 px?}{
    La regla, fijada de antemano, dice que diferencias menores que 0.02 son un empate y que entonces gana el modelo más barato de servir. Los tres mejores empatan, y EfficientNetV2-S a 224 px cuesta 2.9 GFLOPs contra 13.1 de ConvNeXt-T.}
  \qa{28. ¿De dónde sale el umbral de ruido de 0.02?}{
    Es empírico: dos corridas con la misma semilla y la misma configuración difirieron 0.016 de AUC-ROC. Una diferencia menor que eso no distingue un modelo de otro.}
  \qa{29. ¿Por qué se usó la semilla 0 y no la mejor de las tres?}{
    Elegir la mejor semilla sería seleccionar sobre validación y sesgar el resultado. La semilla se fijó antes de ver los números.}
\end{frame}

\begin{frame}{Entrenamiento y selección del modelo (láminas 11 y 12)}
  \qa{30. ¿Qué es EfficientNetV2-S?}{
    Una red convolucional diseñada para dar buena precisión con poco cómputo. Se usó preentrenada en ImageNet-21k. La cabeza es simple: promedio global de los mapas de características, \textit{dropout} y una capa lineal con una sola salida, el logit de melanoma.}
  \qa{31. ¿Por qué subir la resolución de 224 a 384 px no ayudó?}{
    Lo que se midió es que la diferencia (+0.004 de AUC-ROC) queda dentro del ruido, con este conjunto de datos y este presupuesto de cómputo. No se investigó la causa. El ganador del reto sí usó resoluciones altas, pero con muchos más datos y modelos.}
  \qa{32. ¿Cómo se sabe que el entrenamiento no falló en silencio?}{
    Hay un detector de colapso: si las probabilidades de validación casi no varían o el AUC-ROC queda en 0.5, la corrida se marca. Ninguna de las doce corridas colapsó. Además, la línea base de metadatos (AUC-ROC 0.62) confirma que la red aprende de la imagen.}
\end{frame}

% ============================================================
\section{Ensamble y calibración}
\begin{frame}{Ensamble (lámina 13)}
  \qa{33. ¿Por qué ayuda promediar tres redes iguales con distinta semilla?}{
    Cada red comete errores algo distintos por su inicialización y el orden de los datos. Al promediar los logits, parte de esos errores se cancela y baja la varianza de la predicción.}
  \qa{34. ¿Por qué no se usaron la aumentación en inferencia ni los metadatos?}{
    Porque se midieron y no aportaron. Los volteos bajaron el AUPRC del ensamble en 0.008. Añadir edad, sexo y sitio anatómico bajó el AUC-ROC en 0.008 y la especificidad a sensibilidad 0.95 en 0.14. La regla exigía una ganancia mínima y no se alcanzó.}
  \qa{35. En la prueba la mejora es de +0.009 con intervalos traslapados. ¿Por qué desplegar el ensamble?}{
    Porque el criterio se escribió antes de abrir la prueba: desplegarlo si no era peor en AUC-ROC ni en AUPRC. Se cumplió. La mejora es modesta y así se reporta; el costo (tres veces el cómputo, unos 100 ms por imagen) es aceptable.}
\end{frame}

\begin{frame}{Calibración y umbral (lámina 14)}
  \qa{36. ¿Qué significa que una probabilidad esté calibrada?}{
    Que el número se puede leer literalmente: de cada 100 lesiones con una puntuación de 8\,\%, cerca de 8 resultan melanoma. Platt ajusta $p = \sigma(a \cdot \text{logit} + b)$ sobre validación; no cambia el orden de las imágenes ni el AUC.}
  \qa{37. ¿Por qué Platt y no solo escalado por temperatura?}{
    La temperatura corrige la escala, pero no el desplazamiento. El peso de 55 desplaza los logits en aproximadamente $\ln 55 \approx 4$. En el modelo único, la temperatura apenas bajó el Brier de 0.154 a 0.142; Platt lo bajó a 0.016.}
  \qa{38. ¿Qué miden el Brier y el ECE?}{
    Brier: el error cuadrático medio entre la probabilidad y lo que ocurrió (0 o 1). ECE: la diferencia promedio entre la probabilidad predicha y la frecuencia observada, por grupos de puntuación. En ambos, menos es mejor.}
\end{frame}

\begin{frame}{Calibración y umbral (lámina 14)}
  \qa{39. ¿Por qué el umbral es 0.39\,\%? Parece demasiado bajo.}{
    Es el corte que da sensibilidad de al menos 0.95 en validación. Parece bajo porque las probabilidades están calibradas a una tasa base de 1.6\,\%: casi todas las lesiones tienen probabilidades menores que 1\,\%.}
  \qa{40. ¿Por qué sensibilidad 0.95 y no otra?}{
    Es una decisión de costos: un melanoma no detectado es más grave que una consulta innecesaria. La alternativa evaluada, $\tau_{90}$, sube la especificidad a 0.56 pero deja pasar 12 melanomas en vez de 2.}
  \qa{41. ¿Por qué el umbral se calcula como una mediana de remuestreos?}{
    Con 88 melanomas en validación, la sensibilidad de 0.95 depende de solo cuatro casos: un único corte sería frágil. La mediana de 2,000 remuestreos por paciente es más estable.}
\end{frame}

% ============================================================
\section{Resultados}
\begin{frame}{Resultados (láminas 15 a 18)}
  \qa{42. ¿Un AUC-ROC de 0.843 es bueno?}{
    Es razonable para las condiciones: tres redes a 224 px, solo ISIC 2020 y GPU gratuita. El ganador del reto obtuvo 0.949 con 18 modelos, datos de tres años y metadatos; un modelo único de referencia, 0.80. Los conjuntos de evaluación no son los mismos, así que la comparación solo ubica.}
  \qa{43. Con un VPP de 2.6\,\%, ¿para qué sirve?}{
    Para priorizar, no para decidir. Deja fuera al 39\,\% de las lesiones con solo 2 melanomas entre 2,030 no referidas. El VPP bajo es consecuencia directa de la prevalencia de 1.6\,\% y de elegir sensibilidad alta. Es una limitación declarada, la primera de la lista.}
  \qa{44. ¿Por qué el AUPRC bajó de validación (0.132) a prueba (0.096) en el modelo único?}{
    Porque todo se eligió sobre validación, así que sus métricas son optimistas. Esa caída es la brecha de generalización, y es la razón de reportar la prueba.}
\end{frame}

\begin{frame}{Resultados (láminas 15 a 18)}
  \qa{45. ¿El modelo es justo con todos los tonos de piel?}{
    No se puede afirmar. En DDI el AUC-ROC fue parecido en los tres grupos (0.59, 0.60, 0.60), pero bajo en todos: un modelo que discrimina poco no tiene margen para discriminar peor en un grupo. Además, cambian a la vez el tipo de imagen y la tarea.}
  \qa{46. ¿Por qué cae tanto el desempeño en DDI?}{
    Por dos cambios simultáneos: de dermatoscopía a fotografía clínica, y de «melanoma contra el resto» a «maligno contra benigno». Tres algoritmos publicados cayeron a 0.56--0.67 sobre las mismas imágenes; este queda en 0.60.}
  \qa{47. ¿Se puede concluir algo de los subgrupos por edad, sexo o sitio anatómico?}{
    No. Es un análisis exploratorio: cada subgrupo tiene entre 7 y 51 melanomas y sus intervalos se traslapan. Solo detectaría diferencias grandes.}
\end{frame}

% ============================================================
\section{Explicabilidad y sistema}
\begin{frame}{Explicabilidad y artefactos (láminas 19 y 20): preguntas de respaldo}
  \qa{¿Qué muestra exactamente el mapa de activación?}{
    Las regiones que empujan la predicción hacia «melanoma»: una suma de los mapas de la última capa convolucional, ponderada por los pesos de la capa final. Tiene $7 \times 7$ celdas, así que señala regiones, no estructuras finas, y no muestra la evidencia en contra.}
  \qa{¿Qué es la prueba de cordura?}{
    Se reinician los pesos de la red y se recalcula el mapa. Si el mapa casi no cambia, no explica al modelo. Aquí la correlación con el mapa original quedó cerca de cero (umbral: menor que 0.3): el mapa sí depende de lo aprendido.}
  \qa{Si el mapa cae fuera de la lesión, ¿el modelo está mal?}{
    Indica que también usa la piel de alrededor: en los melanomas detectados, 22\,\% de la activación cae en la lesión contra 19\,\% por azar. Dos matices: la máscara de la lesión es automática y burda, y este hallazgo motivó el experimento de artefactos.}
\end{frame}

\begin{frame}{Artefactos y sistema (láminas 20 y 21): preguntas de respaldo}
  \qa{¿Qué demostró el experimento de artefactos?}{
    Que una marca de tinta fuera de la lesión hace referir al 84\,\% de las benignas que no se referían, contra 71\,\% con ruido de control. Y que el ruido hace perder 14 de 84 melanomas. La calidad de la imagen cambia la decisión.}
  \qa{¿Por qué la API usa ONNX y no PyTorch?}{
    Por tamaño y portabilidad: PyTorch instalado ocupa 741 MB; la imagen completa de la API mide 437 MB. La equivalencia con PyTorch se verificó: diferencia máxima de $2.2 \times 10^{-5}$ en las salidas.}
  \qa{¿Por qué la interfaz nunca dice «benigna»? ¿Se guardan las imágenes?}{
    Porque el modelo solo estima melanoma: una probabilidad baja no descarta otros cánceres, y 2 de 86 melanomas no se detectaron. Las imágenes se procesan en memoria y no se almacenan.}
\end{frame}

% ============================================================
\begin{frame}{Preguntas que debe preparar con sus propias palabras}
  \begin{alertblock}{Estas no tienen respuesta escrita aquí a propósito}
    Son personales y el sínodo nota cuando una respuesta es de otra persona.
  \end{alertblock}
  \vspace{0.4em}
  \begin{enumerate}
    \item \textbf{¿Por qué eligió este tema?}
    \item \textbf{¿Qué fue lo más difícil y cómo lo resolvió?}
    \item \textbf{¿Qué haría distinto si empezara de nuevo?}
    \item \textbf{¿Qué aprendió que no sabía al empezar?}
    \item \textbf{¿Usó herramientas de inteligencia artificial? ¿Cómo?}
  \end{enumerate}
  \vspace{0.4em}
  \begin{block}{Sobre la última}
    \footnotesize Conviene revisar antes las reglas de la Facultad y contestar con la verdad, sin rodeos: qué herramientas se usaron, para qué, y qué decisiones y verificaciones fueron propias. Poder explicar cada lámina sin leerla es la mejor respuesta.
  \end{block}
\end{frame}

\end{document}
